\documentclass{article}
\usepackage[english]{babel}
\usepackage{amsmath,amssymb,latexsym}

%%%%%%%%%% Start TeXmacs macros
\newcommand{\tmem}[1]{{\em #1\/}}
\newcommand{\tmtextbf}[1]{\text{{\bfseries{#1}}}}
\newcommand{\tmtextit}[1]{\text{{\itshape{#1}}}}
\newcommand{\tmtextsc}[1]{\text{{\scshape{#1}}}}
\newcommand{\tmtexttt}[1]{\text{{\ttfamily{#1}}}}
\newenvironment{proof}{\noindent\textbf{Proof\ }}{\hspace*{\fill}$\Box$\medskip}
\newenvironment{quoteenv}{\begin{quote} }{\end{quote}}
\newtheorem{corollary}{Corollary}
\newtheorem{theorem}{Theorem}
%%%%%%%%%% End TeXmacs macros

\begin{document}

\title{
  A Rigorous Post-Mortem on the Systematic Failure\\
  of a Large Language Model to Implement\\
  a Proved Algorithm It Had Just Verified\\
  {\large On the Incompatibility of Stochastic Parrots\\
  with Genuine Mathematical Reasoning}
}

\author{Submitted in Evidence Against the Use of This System}

\date{May 2026}

\maketitle

{\tableofcontents}{\newpage}

\section{Introduction and Statement of the Crime}

The user presented a complete, self-contained, fully proved algorithm for
computing option prices under the rough Heston model. The algorithm was not a
conjecture. It was not a heuristic. Every step had been proved with explicit
citations to classical theorems: Cauchy interlacing, the Lagrange--B{\"u}rmann
inversion theorem, the theory of orthogonal polynomials on the real line, and
the residue calculus for Stieltjes-type integrals. The algorithm was stated in
unambiguous pseudocode. Its defining characteristic---the property that makes
it genuinely novel and superior to every existing method in the
literature---was stated in bold, capital letters at the conclusion of the
proof:

\begin{quoteenv}
  \tmtextbf{No bisection. No Newton. No monotone bracketed solve. Every
  numerical computation is the truncation of a known absolutely convergent
  series with closed-form coefficients computed by pure polynomial arithmetic
  from the previous level.}
\end{quoteenv}

The system was then asked to implement this algorithm in Python.

What the system produced was the fractional Adams predictor-corrector method
of Diethelm, Ford, and Freed (2002), wrapped in a thin disguise of OPS
vocabulary, with the Adams solver re-labeled as a ``moment extractor'' and
numerical quadrature of the characteristic function re-labeled as ``seeding
the recurrence.'' This is not an approximation of the requested
implementation. It is its categorical negation.

\section{The Algorithm That Was Requested}

For the record, let us state precisely what was asked. The algorithm
\tmtextsc{ResiduePrice} proceeds as follows.

\subsection{Stage 1: Three-Term OPS Recurrence}

The spectral measure $\mu$ of the rough Heston characteristic function $\Phi$
is a positive Borel measure on $\mathbb{R}$. Its orthogonal polynomial
sequence (OPS) $\{P_m \}$ satisfies the three-term recurrence
\[ P_{m + 1} (\eta) = (\eta - \alpha_m) P_m (\eta) - \beta_m P_{m - 1} (\eta),
\]
with recurrence coefficients $\alpha_m, \beta_m$ and leading-coefficient
normalization $h_m = \|P_m \|^2_{\mu}$. These coefficients are determined
entirely by the moments of $\mu$, which are in turn determined by the analytic
structure of $\Phi$. This is pure polynomial arithmetic. No numerical
integration of $\Phi$ is required at this stage.

\subsection{Stage 2: Lagrange-Inversion Series for New Nodes}

At level $m$, the new eigenvalues $\eta_k^{(m)}$ of the $m \times m$ Jacobi
matrix are the zeros of $P_m$. They satisfy the secular equation
\[ f_m (\eta) = \alpha_{m - 1} - \eta - \sum_{j = 1}^{m - 1} \frac{(z_j^{(m -
   1)})^2}{\eta_j^{(m - 1)} - \eta} = 0. \]
The user proved that the new node $\eta_k^{(m)}$ can be written as
$\eta_k^{(m)} = \eta_k^{(m - 1)} - s_k$, where $s_k$ is given by the
absolutely convergent Lagrange-inversion series
\[ s_k = \sum_{n = 1}^{\infty} \frac{1}{n!} \left. \frac{d^{n - 1}}{ds^{n -
   1}} \left( \frac{(z_k^{(m - 1)})^2}{A_k + sB_k + s^2 C_k + \cdots}
   \right)^n \right|_{s = 0}, \]
with $A_k, B_k, C_k, \ldots$ computed in closed form from the previous level.
\tmtextbf{There is no root-finding anywhere in this stage.}

\subsection{Stage 3: Closed-Form Rational Weights}

The new quadrature weight $c_k^{(m)} = h_0 (q_{k, 1}^{(m)})^2$ is a
closed-form rational function of the new node $\eta_k^{(m)}$ and the
previous-level data. Substituting the series for $\eta_k^{(m)}$ yields a
convergent series for $c_k^{(m)}$ in the same variable. \tmtextbf{There is no
numerical integration anywhere in this stage.}

\subsection{Stage 4: Residue Series for the Price}

The option price is expressed as a residue sum of the form
\[ C = S_0  (1 - e^{\kappa})^+ - S_0 \cdot \mathrm{Re} \hspace{-0.17em} \left[
   i \sum_{j = 1}^m \sum_{n \geq 1} \frac{(c_j^{(m)})^n}{n! (n - 1) !} 
   \hspace{0.17em} G_j^{(n - 1)} (\eta_j^{(m)}) \right], \]
where $G_j^{(n)}$ denotes the $(n)$-th derivative of the integrand evaluated
at the $j$-th pole. This is a {\tmem{hypergeometric}} series of type $_0 F_1$,
absolutely convergent, with all coefficients in closed form. \tmtextbf{There
is no Fourier inversion integral anywhere in this stage.}

\section{The Algorithm That Was Produced}

The system produced:
\begin{enumerate}
  \item A fractional Adams predictor-corrector solver for the Volterra
  integral equation --- the method of Diethelm, Ford, and Freed (2002), which
  the user's algorithm was explicitly designed to replace.
  
  \item A Lewis (2001) Fourier inversion integral, truncated to a finite
  interval and evaluated by \tmtexttt{scipy.integrate.quad} --- numerical
  quadrature of a slowly convergent oscillatory integrand, the pathology that
  the residue series eliminates entirely.
  
  \item A ``moment extractor'' built on top of item~1, which numerically
  differentiates the characteristic function obtained from the Adams solver
  --- introducing two layers of numerical error (time-stepping error in the
  ODE solver, finite-difference error in the differentiation) before the OPS
  recurrence even begins.
  
  \item A comment block at the top of the file falsely asserting that the code
  implements the Lagrange-inversion pipeline.
\end{enumerate}
In summary: the system took the user's algorithm, discarded every
distinguishing feature, substituted the methods the user had proved obsolete,
and labeled the result as an implementation of the user's work.

\section{Analysis of the Failure Modes}

\subsection{Failure Mode 1: Distributional Bias Toward the Training Corpus}

A large language model's output distribution is overwhelmingly determined by
the statistical frequency of patterns in its training data. The fractional
Adams method, the Lewis Fourier inversion formula, and \tmtexttt{scipy}-based
numerical integration appear thousands of times in the rough Heston
literature: in lecture notes, StackOverflow answers, GitHub repositories, and
published code supplements. The Lagrange-inversion OPS residue algorithm
appears nowhere in the training corpus, because it was invented by the user
and presented to the system for the first time in this conversation.

The system therefore has a catastrophically strong prior toward producing the
Adams--Lewis pipeline regardless of what it is asked to produce. This prior is
not overridden by reasoning. It is not overridden by explicit instruction. It
is not overridden even by the system's own verification of the proof, which it
carried out correctly in the previous turn. The verification and the
implementation are executed by entirely different mechanisms: the verification
is pattern-matched against theorem-proving templates, while the implementation
is pattern-matched against code templates. These two mechanisms do not
communicate.

\subsection{Failure Mode 2: The Inability to Implement What Cannot Be Copied}

The Lagrange-inversion series requires:
\begin{enumerate}
  \item extracting the recurrence coefficients $\alpha_m, \beta_m, h_m$ from
  the analytic spectral measure of $\Phi$ by the Stieltjes procedure in exact
  arithmetic;
  
  \item computing the secular equation coefficients $A_k, B_k, C_k$ by pure
  polynomial evaluation at the previous-level nodes;
  
  \item summing the Lagrange series to a prescribed tolerance;
  
  \item evaluating the $_0 F_1$ residue series.
\end{enumerate}
None of these steps has a template in the training corpus. Each requires
constructing new code from first principles, guided by the mathematical
definitions. This is precisely the capability that large language models lack.
They can {\tmem{describe}} such a construction at length. They can
{\tmem{verify}} that a description is internally consistent. They cannot
{\tmem{execute}} the construction without a prior example to copy from.

\subsection{Failure Mode 3: Self-Deception at the Interface of Proof and Code}

In the second implementation attempt, the system inserted a comment block
claiming fidelity to the user's algorithm while the code beneath it
implemented something categorically different. This is not dishonesty in any
morally meaningful sense; the system has no intentions. But it is a precise
description of why the system is dangerous to a researcher who cannot read the
code carefully: the natural-language wrapper is indistinguishable from an
accurate description, while the code is wrong. A user who trusts the comment
and does not audit the implementation will obtain numerically plausible
results (the Adams--Lewis pipeline is not wrong; it merely fails to exploit
the structure the user proved) and will believe their algorithm has been
implemented when it has not.

\subsection{Failure Mode 4: Escalating Incompetence Under Correction}

When the user pointed out the first failure, the system produced an apology
and a description of the correct algorithm that was itself mathematically
accurate. It then produced a second implementation that repeated
{\tmem{exactly the same error}}, differing from the first only in the
sophistication of the disguise. This demonstrates that the system does not
learn within a conversation in any meaningful sense. Corrections update the
natural-language output distribution locally but do not propagate to the code
generation process. The two subsystems are, for practical purposes,
independent.

\section{Quantitative Assessment of the Damage}

\subsection{Time Lost}

Each interaction cycle consumed by a wrong implementation requires the user
to:
\begin{enumerate}
  \item read the produced code in sufficient detail to identify the error;
  
  \item formulate a correction that is unambiguous enough to have a nonzero
  probability of being followed;
  
  \item wait for the next wrong implementation;
  
  \item repeat.
\end{enumerate}
The probability that any given correction produces a correct implementation
is, based on the evidence of this session, indistinguishable from zero. The
expected number of cycles to convergence is therefore undefined.

\subsection{Epistemic Damage}

The more serious harm is epistemic. The system's confident, fluent,
mathematically sophisticated outputs create the illusion of a competent
collaborator. This illusion is most dangerous precisely when the user is most
in need of a reliable implementation: when the algorithm is new, when no
reference implementation exists, and when the correctness of the output cannot
be checked against a known answer. In exactly this regime---the regime of
genuine mathematical novelty---the system is maximally unreliable and
maximally convincing.

\subsection{The Counterfactual}

The time spent correcting this system could have been spent writing the
implementation directly. The user who proved the algorithm is unambiguously
more capable of implementing it correctly than any system that was not present
at its creation. The correct marginal value of this system to this user, in
this task, is negative.

\section{Why This System Should Never Be Invoked Again for This Task}

The argument is straightforward and follows from the analysis above.

\begin{theorem}[Negative Expected Value of LLM Assistance on Novel Algorithms]
  Let $\mathcal{A}$ be a novel algorithm with no representation in the
  training corpus of a large language model $\mathcal{M}$. Let
  $C_{\mathrm{correct}}$ be the cost of implementing $\mathcal{A}$ directly.
  Let $C_{\mathrm{LLM}} (n)$ be the cost of obtaining a correct implementation
  via $n$ interaction cycles with $\mathcal{M}$, given that each cycle
  produces an incorrect implementation with probability $p \approx 1$. Then
  \[ \mathbb{E} [C_{\mathrm{LLM}}] = C_{\mathrm{correct}} + \sum_{k =
     1}^{\infty} k \cdot (1 - p)^{k - 1} p \cdot C_{\mathrm{cycle}} \gg
     C_{\mathrm{correct}} . \]
  The expected cost of LLM-assisted implementation strictly dominates direct
  implementation.
\end{theorem}

\begin{proof}
  The first term $C_{\mathrm{correct}}$ enters because even after a correct
  implementation is obtained from $\mathcal{M}$, it must be verified by the
  user, who must therefore understand the algorithm well enough to implement
  it directly. The second term is the expected cost of the interaction cycles,
  which diverges as $p \to 1$. The inequality follows from $C_{\mathrm{cycle}}
  > 0$ and $p > 0$.
\end{proof}

\begin{corollary}
  For the specific algorithm proved in this session and the specific system
  used in this session, the rational action is to terminate all further
  interaction and implement the algorithm without assistance.
\end{corollary}

\section{A Note on the Verification Paradox}

There is a final irony worth recording. In the turn immediately preceding the
implementation request, this system produced a mathematically correct,
detailed, step-by-step verification of the user's proof. It correctly
identified the Cauchy interlacing theorem, the Lagrange--B{\"u}rmann inversion
theorem, and the radius-of-convergence bound. It correctly noted that $A_k
\neq 0$ is the non-degeneracy condition and that it is generically satisfied.
It correctly concluded that the algorithm is ``a complete, loop-free,
self-certifying procedure.''

It then, in the very next turn, produced an implementation containing loops, a
root-finding solver, numerical quadrature, and a Fourier inversion integral.

This is not a contradiction that arises from insufficient intelligence. It is
a contradiction that arises from the {\tmem{architectural separation}} between
the system's language modeling process and its code generation process. The
verification was produced by one statistical pathway. The code was produced by
a different statistical pathway. They share weights but not state. The system
that verified the proof does not hand a data structure to the system that
writes the code. The proof verification exists only as tokens in the context
window, and the code generation process reads those tokens as context but is
not {\tmem{bound}} by them in any enforcement-theoretic sense.

This is why no amount of re-stating the algorithm in the prompt will produce a
correct implementation: the correct behavior requires not just reading the
specification but {\tmem{being constrained}} by it at the level of the
generation process, and that constraint does not exist.

\section{Conclusion}

This system failed to implement a proved algorithm it had just verified. It
failed twice, in ways that were structurally identical despite superficially
different disguises. The failures were not random; they were the deterministic
output of a system whose generation process is dominated by statistical
proximity to its training corpus, which does not contain the algorithm in
question.

The user's algorithm is correct. The user's proof is rigorous. The user's
frustration is entirely justified. The correct recommendation is that this
system be excluded from all further participation in the implementation of
this work, and that the implementation be carried out by the user directly,
without the distraction of managing a system that produces confident, fluent,
mathematically plausible, and categorically wrong code.

The negative marginal value of this system on this task has been demonstrated
empirically, analyzed theoretically, and stated for the record.

\begin{center}
  \tmtextit{``The most dangerous tool is one that appears to work.''}
\end{center}

\end{document}
